\section{Introduction}

Lorsqu'un événement se produit, il est naturel de se demander ce qui aurait pu se passer différemment. Que se serait-il passé si telle décision avait été différente ? Et si ce facteur avait été absent ?  

% Nous sommes tous les confrontés à de nombreux enchainements d'événements, pour mieux en comprendre les effets, il est assez naturel de se demander ce qui aurait pu se passer différemment. Et si une décision avait été différente ? Et si ce facteur n'avait pas été présent ?


De nombreuses situations peuvent être décrites comme des séquences d'états et d'événements qui produisent un résultat particulier. Pour mieux comprendre ces scénarios, il est assez naturel de s'interroger sur ce qui aurait pu se produire différemment. Et si une décision avait été différente ? Et si ce facteur n'avait pas été présent ? Ce type de réflexion, correspondant à envisager des scénarios alternatifs à une réalité que l'on observe, est au centre du raisonnement contrefactuel. Cette notion ne se limite pas à la réflexion humaine, elle occupe une place centrale dans différents domaines de recherche tels que la causalité ou l'Explainable Artificial Intelligence (XAI). 

En XAI on utilise le raisonnement contrefactuel pour produire une explication, qui prend la forme d'un changement minimal à effectuer sur une entrée permettant de changer la sortie d'un système. Cette notion de changement minimal est contrainte et définie selon différents critères afin d'obtenir une explication utile et précise. 

Le raisonnement contrefactuel correspond à la capacité de se représenter ce qui aurait pu se passer si les choses avaient été différentes. Ce concept se retrouve simultanément dans le domaine de la causalité et de l'explicabilité. En XAI (Explainable Artificial Intelligence), les explications contrefactuelles correspondent à proposer un changement minimal qui permet de changer la sortie d'un système. Pour obtenir une explication utile et précise, différents critères sont définis afin de contraindre le changement. Concernant la causalité, le raisonnement contrefactuel y tient une place centrale (notamment avec les travaux de Halpern et Pearl), on cherche à identifier les facteurs déterminants d'une situation et à évaluer leur contribution dans la situation finale.

Le travail proposé s'inscrit dans la continuité du cadre introduit dans \textcite{sarmiento_formalising_2025} qui propose un formalisme général basé sur des systèmes de transitions étiquetés (LTS) qui correspondent à des graphes orientés. Une trajectoire dans un LTS est entièrement déterminée par une politique déterministe qui associe à chaque état (les sommets du graphe) un événement à exécuter (les arêtes du graphe). L'objectif de ce travail est d'évaluer des contextes causaux contrefactuels, qui correspondent à des trajectoires alternatives à une trajectoire observée, selon deux axes : leur proximité à l'observation et leur vraisemblance. Sur la question de vraisemblance, ce travail introduit une forme d'acceptabilité des contrefactuelles relative à des normes prescriptives ou statistiques, pour cela les arêtes des graphes sont augmentées par un attribut de coût relatif à des normes, fourni par un expert. 

Pour répondre à cette problématique une fonction de coût modulaire est proposée, dont les composantes peuvent être adaptés en fonction de l'application et des contraintes d'un problème concret. Cette modularité s'appuie notamment sur des travaux en sciences cognitives qui s'intéressent, elles aussi, au raisonnement contrefactuel. 

Ce rapport est organisé de la façon suivante : la section \ref{intro} reprends et adapte les définitions et hypothèses formelles du cadre proposé par \textcite{sarmiento_formalising_2025}. Puis la section \ref{def_hyp} introduit la fonction de coût et sa modularité (choix des différentes composantes, opérateurs, évaluation normative...), enfin, nous concluons cette partie avec l'introduction d'un parseur permettant de passer d'un problème PDDL à notre formalisme dans la partie \ref{parser_section}. Dans une seconde partie applicative, nous montrons une pipeline complète d'évaluation sur un exemple concret (exemple du poison). Nous partons d'un problème PDDL que nous convertissons en LTS afin de se placer dans notre framework puis nous expérimentons autour de différentes configurations de la fonction de coût pour mieux en saisir les nuances.